\begin{summarybox}
	\textbf{\textcolor{CSummary}{一句话核心}}
	\textbf{椭圆绕 \(x\) 轴旋转一周所得椭球体积为 \(V=\frac{4}{3}\pi a b^{2}\)，系数 \(\frac{4}{3}\pi\) 与球体相同，只是半径被两半轴“分工”。}

	\medskip
	\textbf{\textcolor{CSummary}{使用条件}}
	\begin{itemize}[leftmargin=*, itemsep=0.5em, topsep=0.4em]
		\item \textbf{条件 1（标准方程）：} 椭圆方程必须为 \(\frac{x^{2}}{a^{2}}+\frac{y^{2}}{b^{2}}=1\) 的形式，且 \(a>0,b>0\)。
		\item \textbf{条件 2（明确旋转轴）：} 旋转轴为 \(x\) 轴，此时截面半径由 \(y\) 坐标确定，公式中 \(b\) 带平方。
	\end{itemize}

	\medskip
	\textbf{\textcolor{CSummary}{关键提醒}}
	\begin{itemize}[leftmargin=*, itemsep=0.5em, topsep=0.4em]
		\item \textbf{易错点（半轴与轴长）：} \(a,b\) 是半轴长（即轴长的一半），切勿将全长直接代入。
		\item \textbf{检查项（先化标准）：} 若方程非标准形式，务必先变形为 \(\frac{x^{2}}{a^{2}}+\frac{y^{2}}{b^{2}}=1\)，再读取 \(a,b\)。
	\end{itemize}
	\begin{center}
		% fig08_summary_flow_clean_v6.tex
% 用法：主文件提前加载：
% \usepackage{tikz,amsmath,graphicx}
% \usetikzlibrary{arrows.meta,positioning,calc}

\begin{center}
	\begin{tikzpicture}[
			x=1cm,
			y=1cm,
			>=Stealth,
			line cap=round,
			line join=round,
			every node/.style={font=\small},
			flowline/.style={
					line width=0.9pt,
					blue!70!black
				},
			arrow/.style={
					->,
					line width=0.9pt,
					blue!70!black,
					shorten >=2pt,
					shorten <=2pt
				},
			topbox/.style={
					draw=blue!65!black,
					fill=blue!3,
					rounded corners=7pt,
					line width=0.95pt,
					align=center,
					minimum height=1.35cm,
					inner xsep=8pt,
					inner ysep=7pt
				},
			formulaBox/.style={
					draw=blue!65!black,
					fill=blue!3,
					rounded corners=7pt,
					line width=0.95pt,
					align=center,
					minimum width=3.55cm,
					minimum height=1.65cm,
					inner xsep=8pt,
					inner ysep=7pt
				},
			readBox/.style={
					draw=blue!65!black,
					fill=blue!3,
					rounded corners=7pt,
					line width=0.95pt,
					align=center,
					minimum width=3.55cm,
					minimum height=1.95cm,
					inner xsep=8pt,
					inner ysep=7pt
				},
			judgeBox/.style={
					draw=red!55!black,
					fill=red!3,
					rounded corners=7pt,
					line width=0.95pt,
					align=center,
					minimum width=3.55cm,
					minimum height=0.95cm,
					inner xsep=8pt,
					inner ysep=6pt,
					font=\small\bfseries
				},
			resultX/.style={
					draw=green!50!black,
					fill=green!4,
					rounded corners=7pt,
					line width=1pt,
					align=center,
					text width=4.65cm,
					minimum height=2.65cm,
					inner xsep=10pt,
					inner ysep=9pt
				},
			resultY/.style={
					draw=orange!70!black,
					fill=orange!5,
					rounded corners=7pt,
					line width=1pt,
					align=center,
					text width=4.65cm,
					minimum height=2.65cm,
					inner xsep=10pt,
					inner ysep=9pt
				},
			noteBox/.style={
					draw=red!55!black,
					fill=red!3,
					rounded corners=7pt,
					line width=0.95pt,
					align=center,
					text width=11.15cm,
					minimum height=1.25cm,
					inner xsep=12pt,
					inner ysep=8pt
				},
			arrowLabel/.style={
					font=\scriptsize,
					fill=white,
					inner xsep=3pt,
					inner ysep=1pt,
					text=black!65
				}
		]

		% ===================== 标题 =====================
		\node[
			font=\Large\bfseries,
			text=blue!60!black
		] at (6.2,7.05)
		{一图流总结：先标准，后读轴，再看旋转轴};

		% ===================== 第一层：标准化与读轴 =====================
		\node[topbox, minimum width=3.1cm] (origin) at (1.55,5.55)
		{原方程\\[4pt]
		{\normalsize 可能不是标准式}};

		\node[formulaBox] (standard) at (6.2,5.55)
		{化为标准式\\[3pt]
			\(\displaystyle \frac{x^2}{a^2}+\frac{y^2}{b^2}=1\)};

		\node[readBox] (read) at (10.85,5.55)
		{读半轴\\[4pt]
			\(\displaystyle a:\ x\text{ 方向}\)\\[3pt]
			\(\displaystyle b:\ y\text{ 方向}\)};

		\draw[arrow]
		(origin.east) -- node[arrowLabel, above] {右边化为 \(1\)}
		(standard.west);

		\draw[arrow]
		(standard.east) -- node[arrowLabel, above] {看分母}
		(read.west);

		% ===================== 第二层：判断旋转轴 =====================
		\node[judgeBox] (judge) at (6.2,3.70)
		{判断旋转轴};

		\draw[arrow] (standard.south) -- (judge.north);

		% ===================== 第三层：两个结论卡片 =====================
		% 关键：用 anchor=north 固定卡片上边界。
		% 这样即使中文字体导致卡片变高，也只会向下扩展，不会向上顶到箭头。

		\node[resultX, anchor=north] (xaxis) at (3.15,2.20)
		{{\bfseries 绕 \(x\) 轴}\\[5pt]
		截面半径来自 \(y\) 方向\\[7pt]
		\(\displaystyle V_x=\frac{4}{3}\pi ab^2\)\\[4pt]
		{\scriptsize \(x\) 方向半轴 \(a\) 只出现一次}};

		\node[resultY, anchor=north] (yaxis) at (9.25,2.20)
		{{\bfseries 绕 \(y\) 轴}\\[5pt]
		截面半径来自 \(x\) 方向\\[7pt]
		\(\displaystyle V_y=\frac{4}{3}\pi a^2b\)\\[4pt]
		{\scriptsize \(y\) 方向半轴 \(b\) 只出现一次}};

		% 分叉线：不要再在线上写“绕 x 轴 / 绕 y 轴”，避免遮挡。
		\coordinate (split)     at (6.2,2.72);
		\coordinate (leftDrop)  at (3.15,2.72);
		\coordinate (rightDrop) at (9.25,2.72);

		\draw[arrow]    (judge.south) -- (split);
		\draw[flowline] (leftDrop) -- (rightDrop);
		\draw[arrow]    (leftDrop) -- (xaxis.north);
		\draw[arrow]    (rightDrop) -- (yaxis.north);

		% ===================== 第四层：口诀 =====================
		\node[noteBox, anchor=north] (note) at (6.2,-1.25)
		{{\bfseries 口诀：}
			绕哪条轴，哪条轴方向的半轴只出现一次；\quad
			垂直于旋转轴的方向作为圆半径，进入面积时要平方。};

		% ===================== 图注 =====================
		\node[
			font=\small,
			text=black!65,
			anchor=north
		] at (6.2,-2.75)
		{图：一图流总结卡};

	\end{tikzpicture}
\end{center}
	\end{center}
	\medskip
	\textbf{\textcolor{CSummary}{新增两种证明的速记卡}}
	\begin{itemize}[leftmargin=*, itemsep=0.5em, topsep=0.4em]
		\item \textbf{沿 \(x\) 切（原有圆盘法）：}截面是圆，\(S(x)=\pi b^{2}(1-x^{2}/a^{2})\)，积分范围 \([-a,a]\)。
		\item \textbf{沿 \(y\) 切（新增椭圆截面法）：}旋转轴\textbf{仍为 \(x\)}；固定 \(y\)，截面在 \(xz\) 平面，\(S(y)=\pi ab(1-y^{2}/b^{2})\)，积分范围 \([-b,b]\)。
		\item \textbf{球体伸缩（新增方程法）：}球 \(u^{2}+y^{2}+z^{2}=b^{2}\)，令 \(x=(a/b)u\)，得到椭球方程；沿 \(x\) 的厚度乘 \(a/b\)，因此 \(V=(a/b)\cdot\frac43\pi b^{3}\)。
	\end{itemize}
	\[
	\boxed{\int_{-a}^{a}\pi b^{2}\left(1-\frac{x^{2}}{a^{2}}\right)dx
	=\int_{-b}^{b}\pi ab\left(1-\frac{y^{2}}{b^{2}}\right)dy
	=\frac{a}{b}\cdot\frac43\pi b^{3}
	=\frac43\pi ab^{2}}
	\]
	\textbf{记忆推广：}一般椭球的三个半轴为 \(a,b,c\) 时，\(V=\frac43\pi abc\)；本题两个垂直于旋转轴的半轴都为 \(b\)。

\end{summarybox}
